\documentclass[output=paper,colorlinks,citecolor=brown]{langscibook}
\ChapterDOI{10.5281/zenodo.18845592}
\author{Saloumeh Gholami\orcid{}\affiliation{Cambridge University} and Ulrich Mehlem\orcid{}\affiliation{Goethe Universität Frankfurt}}
\title{Introduction: Critical issues in migration, language, and identity}
\abstract{This introduction investigates the intricate connections between migration, identity, and language, addressing significant gaps in existing research. It explores the evolution of identity, moving from classical philosophical ideas to modern concepts such as hybrid and transnational identities. Identity is portrayed as an ever-changing construct shaped by both individual and collective experiences, as well as by the broader forces of globalisation and migration.

A key focus is the critique of migration studies' tendency to prioritise dominant groups, often neglecting the nuanced experiences of minority populations within minorities. This oversight \citep{esser2006,knigge_klinger_schnoor_gogolin2015,Welker2021RelativeEducationRefugees} limits the understanding of the diverse challenges these communities face.

The discussion highlights the critical role of migrants' countries of origin in shaping their identities. Factors such as political instability, economic hardships, and linguistic diversity are examined for their influence. Through case studies of multilingual families, the text illustrates how heritage language use contributes to educational outcomes and cultural preservation.

Additionally, the introduction critiques the oversimplification of the MENA region’s rich diversity. It underscores how systemic discrimination and conflict disproportionately affect minorities, leading to forced migration and the erosion of cultural heritage.

The text also delves into the challenges of categorising languages in migration contexts. By advocating for terms like \textit{heritage language}, it captures the fluid and adaptive nature of language use in these settings. Overall, this introduction emphasises the need for more inclusive and nuanced approaches to understanding migration, identity, and language.}

\IfFileExists{../localcommands.tex}{
   \addbibresource{../localbibliography.bib}
   \usepackage{tabularx,multicol}
\usepackage{url}
\urlstyle{same}
\usepackage{textcomp}
 
\usepackage{langsci-optional}
\usepackage{langsci-lgr}
\usepackage{langsci-gb4e}
\usepackage{langsci-branding}

% \usepackage{fontspec}
\babelprovide{georgian} % Georgian language
\babelprovide{armenian} % Armenian language
\babelfont[georgian,armenian]{rm}{DejaVu Sans}
\newfontfamily\georgianfont{DejaVu Sans}
\newfontfamily\armenianfont{DejaVu Sans}

\usepackage{dialogue}
\usepackage{attrib}
\usepackage{attrib} % speaker labels in Koca's chapter
\usepackage{phoenician} % alf, bet in Mengozzi's chapter

   

\makeatletter
\let\thetitle\@title
\let\theauthor\@author
\makeatother

\newcommand{\togglepaper}[1][0]{
   \bibliography{../localbibliography}
   \papernote{\scriptsize\normalfont
     \theauthor.
     \titleTemp.
In Saloumeh Gholami \& Ulrich Mehlem (eds.), \textit{Identity formation, language and migration dynamics: Minorities in the MENA region and the German diaspora}, Berlin: Language Science Press [preliminary page numbering]
   }
   \pagenumbering{roman}
   \setcounter{chapter}{#1}
   \addtocounter{chapter}{-1}
}

\newbool{bookcompile}
\booltrue{bookcompile}
\newcommand{\bookorchapter}[2]{\ifbool{bookcompile}{#1}{#2}}


\newcommand{\georgian}[1]{{\georgianfont #1}}
\newcommand{\armenian}[1]{{\armenianfont #1}}

\newcommand{\textarab}[1]{{\arabicfont \beginR #1} \endR}
\newcommand{\texthebrew}[1]{{\hebrewfont \beginR #1} \endR}
\newcommand{\textsyriac}[1]{{\syriacfont \beginR #1} \endR}


\newfontfamily\phoenicianfont{Tuffy Regular.ttf}

   %% hyphenation points for line breaks
%% Normally, automatic hyphenation in LaTeX is very good
%% If a word is mis-hyphenated, add it to this file
%%
%% add information to TeX file before \begin{document} with:
%% %% hyphenation points for line breaks
%% Normally, automatic hyphenation in LaTeX is very good
%% If a word is mis-hyphenated, add it to this file
%%
%% add information to TeX file before \begin{document} with:
%% %% hyphenation points for line breaks
%% Normally, automatic hyphenation in LaTeX is very good
%% If a word is mis-hyphenated, add it to this file
%%
%% add information to TeX file before \begin{document} with:
%% \include{localhyphenation}
\hyphenation{
    par-a-digm
    Ver-to-vec
    Go-go-lin
    sa-gen
    mo-misch
    ant-wor-te-te
    dix-wa-zim
    spre-chen
    B-Verf-GE
    Swe-dish
    re-mains
    schrei-ben
    di-pey-vin
    Chris-ten-sen
    Wes-sen-dorf
    Ja-va-khi-shvi-li
    ana-lysis
}

\hyphenation{
    par-a-digm
    Ver-to-vec
    Go-go-lin
    sa-gen
    mo-misch
    ant-wor-te-te
    dix-wa-zim
    spre-chen
    B-Verf-GE
    Swe-dish
    re-mains
    schrei-ben
    di-pey-vin
    Chris-ten-sen
    Wes-sen-dorf
    Ja-va-khi-shvi-li
    ana-lysis
}

\hyphenation{
    par-a-digm
    Ver-to-vec
    Go-go-lin
    sa-gen
    mo-misch
    ant-wor-te-te
    dix-wa-zim
    spre-chen
    B-Verf-GE
    Swe-dish
    re-mains
    schrei-ben
    di-pey-vin
    Chris-ten-sen
    Wes-sen-dorf
    Ja-va-khi-shvi-li
    ana-lysis
}

   \boolfalse{bookcompile}
   \togglepaper[1]%%chapternumber
}{}

\begin{document}
\maketitle

\isi{Migration} remains a critical and multifaceted issue that has garnered extensive scholarly attention, particularly in recent years. New developments and ongoing conflicts across the Middle East and North Africa (MENA) region have driven successive waves of migration, with many migrants seeking refuge in Europe. These movements have intensified research interest, resulting in a substantial body of academic work exploring the diverse dimensions of migration.
\largerpage

However, despite this extensive body of research, substantial limitations persist in migration studies, largely owing to the inherent complexity of migration as a phenomenon. \isi{Migration} is shaped by a nuanced and interwoven mix of economic, political, social, and cultural factors, creating interactions that resist straightforward categorisation and analysis. This complexity poses a challenge for current research to fully encapsulate the diverse experiences of migrants, particularly those of \isi{minority} groups, whose perspectives are frequently underrepresented in existing literature.

Central to these challenges is the conceptual ambiguity surrounding ``identity'', a term often applied so broadly that it risks losing analytical precision. To grasp its complexity, it is essential to trace the evolution of identity from its classical philosophical origins through to its reconceptualisations in modern social sciences. Understanding how identity has been shaped and reshaped in response to historical, social, and technological shifts -- particularly within the contexts of migration and language -- is critical to deepening the analysis of migrant experiences and the transformative impact of migration on individual and group identities.

\section{The evolving complexity of identity: Historical challenges and theoretical layers}\label{Intro.ss.1}

In the classical era, especially within Hellenic philosophy, identity was a foundational topic, with thinkers analysing its essential properties to define both commonalities and distinctions. This approach hinged on the idea of \textit{differentia specifica}, a method of categorising entities by their shared and unique traits. Accordingly, two entities (A and B) were deemed identical only if they shared all properties without distinguishing differences; in such cases, their separate labels (A, B) represented mental constructs rather than inherent differences \citep{windelband1910}. This line of thought by Socrates, Plato, and Aristotle established a foundation for Western notions of personal identity, focusing on intrinsic virtues and characteristics that differentiate individuals.
\largerpage

During the medieval period, at least within Europe, identity was framed by religious affiliation and collective associations such as estate, profession, and local community, rather than individual traits. The Renaissance revived interest in personal identity, inspired by classical ideals of human potential and individuality, while the Enlightenment expanded this view to include the individual's role within the emerging nation-state. Social class also became a key identifier, reflecting society's broader structures. In modern times, identity encompassed not only social roles \citep{tajfel_billig_bundy1971,tajfel_turner1979} but also national affiliations reshaped by industrialisation and globalisation \citep{brubaker_cooper2000}.

In the digital era, the concept of identity has become increasingly complex, shaped by the interplay of real-world interactions and virtual personas. Among the collectives influencing group identities, linguistic, religious, and so-called ``ethnic'' affiliations retain a significant role. The term ``ethnic'' here refers to social groups sharing a common ancestry, language, religion, territory, or cultural practices, distinguishing them from groupings by \isi{gender}, profession, or socio-economic status \citep{peoples_bailey2010}. As Max Weber highlighted, such definitions rest not on objective criteria but on the group's self-perceived unity \citep{weber1968}.

Modern social sciences and psychology distinguish between personal identity and group identity. Personal identity can be understood as a continuous process of constructing a coherent self throughout one's life, forming a ``vertical dimension'' of identity \citep{erikson1980}. In parallel, individuals navigate the ``horizontal dimension'' by balancing roles across different groups -- such as family, profession, religion, and community \citep{goffman1963}. While group affiliations may have less influence today, individuals remain distinct from their cultural or linguistic associations, as Jullien observes: ``Il faudra donc repenser le rapport du sujet à la culture dès lors que celle-ci n'est pas posée comme ``sa'' culture''  `It will therefore be necessary to rethink the relationship of the subject to culture as soon as it is not considered as his/her culture' (\citealt[57]{jullien2016}, \citealt{jullien2021}). Group identities continue to shape \isi{self-perception}, influenced by external perceptions and interactions.

In contemporary discourse, concepts such as ``transnational identity'' \citep{vertovec1999} and ``hybrid identity'' \citep{bhabha1994,hall1990} illustrate how individuals construct identities spanning home and host cultures, particularly in globalised and postcolonial contexts. These frameworks highlight the complex adaptations individuals make when navigating new environments. A notable precursor to such views was the philosopher Hannah Arendt, who noted in \textit{The human condition}: ``We can only learn who somebody is or was, if we hear the story whose hero he himself is, his biography'' \citep{arendt2020,arendt1998}. Poststructuralist and postcolonial theories extend this notion, suggesting that the ``self'' is represented through an endless process of construction rather than reaching any fixed essence. Following Derrida's concept of \textit{différance}, \citet{hall1990} views identity as an ongoing process of positioning the self within social power structures. Similarly, Jullien asserts that culture itself unfolds within a dynamic tension between plurality and unity: ``Je dirai plutôt que le propre du culturel est qu'il se déploie dans cette tension -- ou cet écart -- du pluriel et de l'unitaire'' `I would rather say that the specificity of culture is that it unfolds in this tension -- or this gap -- of plurality and the unitary' \citep[46]{jullien2016}.

In contemporary discourse, the term ``identity'' is often applied too broadly, leading to oversimplification, or conversely, in an excessively narrow sense that limits its inherent complexity. Such narrow interpretations risk reducing identity to rigid, exclusive categories that fail to reflect its dynamic and multidimensional nature \citep{gholami2023,brubaker_cooper2000}. \isi{Identity}, rather than being a static or immutable construct, comprises multiple dimensions shaped by diverse factors, adapting continuously in response to changing contexts.

A pertinent example illustrating this complexity is that of \ili{Kurdish} Jews who immigrated from Iraqi Kurdistan to Israel in 1950 \citep{shwartz-beeri_yisrael2000}. Their identity weaves together various intersecting dimensions: religious (Jewish), ethnic (\ili{Kurdish}), national (initially Iraqi and later Israeli), and historical (as post-1950 immigrants). Each layer contributes to a unique, composite identity that defies simplistic categorisation, underscoring the dynamic interplay between these overlapping affiliations and the evolving nature of identity.

\section{Emphasis on dominant groups with limited attention to minority populations within minorities}\label{Intro.ss.2}
\largerpage
A critical limitation within migration studies, especially concerning the \isi{MENA region}, is the field's predominant focus on majority ethnic, religious, or linguistic groups. This dominant-group-centric perspective often oversimplifies the diverse identities of \isi{minority} populations within broader migrant communities, leading to a lack of nuanced understanding of \isi{minority} populations within \isi{minorities}. For example, studies of \ili{Iranian} migration often foreground Shiite Muslim identity, which, while central, fails to capture the unique experiences of religious \isi{minorities} such as Jews, Zoroastrians, Christians, and Ahl-e Haqqs \citep{gholami2023}. By neglecting these layered \isi{minority} identities, researchers risk overlooking the complex cultural and social dynamics that these groups navigate, including how they define themselves in contrast to both their original and host communities.

This focus on dominant groups extends into academic frameworks as well, where many sociological theories inherently reflect majority-group perspectives. \citegen{berry1997} acculturation theory, a foundational model in migration studies, categorises migrant integration along two axes (community of origin and host society), yielding four main outcomes: integration, assimilation, ethnic separation, and marginality. While influential, this model primarily reflects majority group contexts, often rendering \isi{minority} identities as subsets of these broader categories without sufficient consideration for their distinctive experiences. Similarly, the Subjective Ethnolinguistic Vitality Questionnaire \citep{bourhis_giles_rosenthal1981} measures sociological factors -- such as status, demography, and institutional support -- within ethnic groups but has often been critiqued for reifying reference categories that may not fully capture \isi{minority} identities within \isi{minorities}. Critics like \citet{husband_khan1982} and \citet{tollefson1991} argue that such frameworks often embed majority perspectives, limiting their relevance to diverse migrant communities.

\isi{Communication Accommodation Theory (CAT)}, which emerged as a model to explain how individuals adjust their communication styles in social contexts, also highlights challenges in recognising the complexities of \isi{minority} identities. CAT distinguishes between convergence (adapting to another's communication style), divergence (emphasising distinctiveness), and maintenance\is{Language!maintenance} (retaining one's communicative style) \citep{Giles2016CommunicationAccommodation}. However, understanding how these dynamics manifest within \isi{minority} groups requires insight into sociohistorical contexts and power imbalances that mainstream applications of CAT may overlook. Consequently, while the model acknowledges social and linguistic adaptation, its application to \isi{minority} identities requires deeper contextualisation, particularly when analysing interactions within and between \isi{minorities}.

This volume aims to address these gaps by advancing migration studies to include not only dominant groups but also the nuanced identities within \isi{minority} populations. By foregrounding multilingualism, cultural networks, and superdiversity -- where migrants may constitute \isi{minorities} even within their homelands -- this work seeks to illuminate the layered and intersectional realities of \isi{minority} identities in migration contexts, offering a more comprehensive understanding of identity and integration beyond conventional frameworks.

\section{The problem of overlooking countries of origin in migration studies and identity formation}\label{Intro.ss.3}

A fundamental problem in the field of migration research is the predominant focus on destination countries, which often overshadows the crucial role played by countries of origin. Factors such as economic conditions, political instability, social dynamics, and cultural traditions in migrants' homelands significantly influence identity formation. However, these factors often receive insufficient scholarly attention, despite their undeniable impact. While extensive research exists on assimilation, acculturation, and the various linguistic shifts occurring in \isi{diaspora} contexts, there remains a conspicuous gap in understanding the rich linguistic and cultural tapestries from which these migrants emerge.

For instance, debates around multilingualism in countries like Germany illustrate this imbalance \citep{gogolin_neumann}. New assimilation theories \citep{alba_nee2003,alba_nee1997} have historically prioritised the acquisition\is{Language!acquisition} of the host country's language as a measure of social and cultural integration. In this view, maintaining other languages is often seen as detrimental, a notion reinforced by studies like \citet{esser2006} and \citet{hopf2005} which base their arguments on the ``time on task'' hypothesis \citep{Carroll1963ModelSchoolLearning}. According to this hypothesis, allocating more time to learning the host country's language results in better language proficiency\is{Language!proficiency}. PISA studies provide further evidence, showing correlations between the use of a non-majority language at home and academic performance, categorising students based on their migration background \citep{mullis_martin_foy_hooper2017}.

By contrast, the theory of segmented assimilation \citep{portes_rumbaut2001} stipulated a continuous importance of migrants' first languages even after long periods of settlement in the immigration country. Again, on the basis of the international PISA data, \citet{agirdag_vanlaar2016} differentiated the analysis of language usage between parents, siblings and peers and showed a positive correlation of L1-usage with the father and educational success. In the same line, \citet{rauch2019} argued that the trend towards less L1-usage of young people with immigration background in the PISA study does not hold for siblings and friends. On the other side, educational success is not significantly related to less L1-usage inside of the family (137). Only the L1-usage with best friend outside of school exhibits such a relationship \citep{rauch2019}.

\citet{knigge_klinger_schnoor_gogolin2015} tested the influence of L1 proficiency\is{Language!proficiency} on L2 proficiency\is{Language!proficiency} of young migrants of Turkish, Russian and Vietnamese background in Germany according to the interdependency hypothesis \citep{cummins1979}. This influence is shown to be small (6\% of variance explained), but significant. In terms of national background, their data showed an advantage in L2 proficiency\is{Language!proficiency} of the Vietnamese, followed by the Russian and the lowest by the \ili{Turkish} speakers, even after control for socio-economic background.

In such a perspective, the longitudinal study of \citet{leseman_scheele_mayo_messer2009} in the Netherlands differentiated heritage language and country of origin. It compared children of \ili{Dutch}, \ili{Turkish} and Moroccan \ili{Berber} origin in their dual language development from preschool to primary school, taking into account the socio-economic and cultural background of the parents and the language learning environment at home. This background is, for the Moroccan \ili{Berber} group, characterised by a very complex linguistic situation. While Tarifit\il{Berber!Tarifit} \ili{Berber} is spoken at home, the official language of \isi{Morocco} is Modern Standard Arabic\il{Arabic!Modern Standard}, which is also present in the media to some extent. Additionally, Moroccan Arabic\il{Arabic!Moroccan} is needed in many situations of everyday life, and even \ili{French} and \ili{Spanish} are present.

In the Netherlands, media use inside the Moroccan families is multilingual, while the \ili{Turkish} group only receives input in two languages. Therefore, the impact on vocabulary development in L1 (Berber) is smaller than for the Turkish children. But even in that group, after controlling for \ili{Dutch} vocabulary knowledge, short-term memory and Raven IQ at age 3, the variance of \ili{Dutch} proficiency\is{Language!proficiency} at the age of six years is also explained to some extent by the oral proficiency\is{Language!proficiency} in the first language at age 3 (albeit with small effect sizes, \citealt{leseman_scheele_mayo_messer2009}). In addition to such studies, it is essential to understand that migration is not just about the destination, but equally about the point of departure. To have a comprehensive grasp of the sociolinguistic implications of migration, it is crucial to delve deeper into the roots of migrants: their hometowns, the languages they speak, the cultural norms they adhere to, and the complexities of their original environments. 

\section{Oversimplifying the complexity of the MENA region}\label{Intro.ss.4}

Another critical issue in migration studies is the tendency to oversimplify the \isi{MENA region}'s inherent complexity. This region is marked by a rich diversity of ethnicities, religions, and cultural practices, each adding layers to its social and political landscape. The coexistence of multiple groups with distinct traditions and beliefs creates a dynamic yet fragile balance of cohabitation, often accompanied by episodes of tension or conflict. A nuanced understanding of these intricacies is essential for any comprehensive analysis of migration from the \isi{MENA region}, as it reveals the specific cultural, religious, and social forces that shape migrants' experiences and the formation of their identities. 

The \isi{MENA region} has been the site of numerous conflicts and wars, which have disproportionately affected \isi{minority} populations. These conflicts often result in significant displacement, loss of life, and destruction of cultural heritage, exacerbating the vulnerabilities of these groups. Understanding the particular cultural, religious, and social forces that influence migrants from this region is essential for any comprehensive analysis.

For instance, the prolonged Syrian Civil War has displaced over 10 million people, finding refuge in Turkey, Lebanon, Jordan and Europe, many of whom belong to \isi{minority} groups like the Kurds, whose experiences of war and displacement are intertwined with the unique socio-political status of regions like Rojava. Similarly, in Northern Iraq, ISIS's genocide against the Êzȋdi community in 2014 reveals how sectarian and ethnic conflicts disproportionately impact vulnerable groups, prompting forced migration and exacerbating cultural erasure \citep{schmidinger2019}. In Turkey, after the Suruç bombing in 2015, armed conflict started again in \ili{Kurdish} areas of East Anatolia after a period of ``\ili{Kurdish} opening'' of the Erdoğan regime. Political protest in Iran in 2022 started when a young woman of \ili{Kurdish} origin, Mahsa Amini, died after having been arrested by the religious morality police \citep{Kazemzadeh2023MassProtestsIran}. 

In Afghanistan, decades of conflict have entrenched divisions among Pashtuns, Tajiks, Hazaras, and Uzbeks, with each wave of regime change impacting these communities differently. The Taliban's recent re-assertion of \ili{Pashto} dominance is but one manifestation of how ethnic dynamics shape migration patterns and individual migrant experiences \citep{ahwar2023}. The constitution of 1964 declared \ili{Pashto} and \ili{Dari} as official, but only \ili{Pashto} as the national language. Other languages such as Uzbek, Turkmen, Baluchi and Nuristani were labelled regional languages. During Soviet military intervention and civil war, ethnic issues grew in importance. The first Taliban regime (1996--2001) lead to the dominance of the Pashtuns \citep{nawid2012}. After the US-led intervention in 2001 and the reconstruction of a Democratic regime, the balance between different groups was reestablished \citep{schetter2005}. When civil war started again, particularly \ili{Pashto}-speaking provinces suffered. 

These dynamics are not restricted to regions with intense conflict but also extend to areas like \isi{Morocco}, where the \il{Berber!Amazigh|see{Berber!Tamazight}} (\ili{Berber}) community in the Rif Mountains has experienced cyclical marginalisation and periods of state repression \citep{kratochwil2002}. Although the integration of the \il{Berber!Amazigh|see{Berber!Tamazight}} language into the curriculum was a positive step, recent unrest underscores the fragility of ethnic relations. 

Economic instability in the region poses another layer of difficulty for \isi{minorities}. These groups often face higher rates of poverty, unemployment, and limited access to essential services. The economic marginalisation of \isi{minorities} not only contributes to their immediate hardship but also has broader implications for their long-term social and political integration.

\isi{Discrimination} against \isi{minorities}, both in legal systems and societal norms, is a pervasive issue in the region. This discrimination can manifest in various forms, including limited access to justice, education, and employment opportunities, and can lead to systemic exclusion and inequality. Examining the legal and social frameworks that perpetuate such discrimination is vital for understanding the challenges faced by \isi{minorities}.

The phenomena of migration and forced displacement have significantly reshaped \isi{minority} communities in the MENA region. These movements not only result in demographic changes but also affect the cultural and social fabric of both the communities left behind and those that receive new members. The transformation of \isi{minority} communities due to migration and displacement is a critical area of study, with implications for cultural preservation and community cohesion. The dynamics of \isi{minority} issues in the \isi{MENA region} are further complicated by global and regional factors. International political decisions, economic sanctions, and regional power struggles can have a cascading effect on \isi{minority} populations, influencing everything from policy decisions to grassroots movements.

In line with the need to avoid oversimplifying the \isi{MENA region}'s intricate sociocultural fabric, Brizić's mixed methods study, The Sound of Inequality \citep[]{brizic2021,brizic_bulut_simsek2021} provides a pertinent example of a migration study that addresses the complex conditions of migrants' countries of origin. 

The study focuses on narrative-biographical interviews with 115 parents of different ethnic origin from Kurdistan/Turkey in Austria which were conducted in \ili{Kurmanci}-\ili{Kurdish}, \ili{Turkish} and \ili{German}.\footnote{%
The second group of migrants from Serbia will not be discussed here.} The speakers belong to a variety of different ethnic groups: 30 families are \ili{Kurdish} speakers from central Anatolia, where \ili{Kurmanci} \ili{Kurdish} is already used as a kind of lingua franca besides Bedeyni-\ili{Kurdish} und Caucasian languages. As a second group, 14 \ili{Zaza}/\ili{Kurdish}-speaking families from east or southeast Anatolia were distinguished. But also, the remaining 71 predominantly \ili{Turkish} speaking families from different regions list a variety of other languages and varieties that played or play some role in their environment. 

The biographical narratives of the parents are analysed as voices in Blommaerts' sense of the term as ``attempts … to make themselves understood by others… the practical conversion of socially loaded resources into socially loaded semiotic action''. \citep[427]{blommaert2008}, leading to a typology of four voices, defined as the Resisting, Targeting, Intermediate, and Balancing Voice. While the resisting voice referred to the destroyed \ili{Kurdish} village as a symbol of group identity and place of \isi{belonging} in the past, against which the present situation of migration and acculturation is rejected, the targeting voice is about leaving this past behind and aspiring education and better living conditions in migration, taking language loss into account. The intermediate voice, on the other side, emphasises education, and favours the official language of the country of origin, \ili{Turkish}, without completely suppressing the \ili{Kurdish} home variety. The actual situation in Austria is perceived as transitional. The balanced voice finally reconciles the variety of origin with a perspective of educational success in the new host country. These four voices correspond to some extent with the socio-economic background of the family and the multilingualism of the community of origin: lower educational background of the mothers corresponds with the resistance and targeting voices, while in half of the balancing voices, the mothers completed at least 10 school years \citep{brizic2021}. With respect to multilingualism, the families of the intermediate and, to a lesser extent, the targeting voices use at least three and up to six languages, while the balanced and resisting voices restrict themselves to two. In this way, the experiences of multilingualism, discrimination and political oppression in the country of origin shape attitudes\is{Language!attitudes}, which are still important when bringing up their children in the immigration country Austria. 

\section{The challenge of accurately categorising and describing languages in migration contexts}\label{Intro.ss.5}

The layered difficulties involved in effectively categorising and describing the multilingual realities of migrant populations underscore the need for careful, flexible research approaches that account for both the social meanings of language -- considering both emic and etic perspectives -- and the fluidity of language use within these communities. In premodern times, different varieties of the same \textit{language} and different languages coexisted in one area under the roof of an Empire with its official language or language of administration. In such a way, Europe during the Middle Ages was multilingual, even inside a territory that was supposed to constitute later France or England. The language of literacy was \ili{Latin}, while local varieties slowly made their way into written usages such as poetry. 

In the \isi{MENA region}, \ili{Arabic}, \ili{Persian} and \ili{Ottoman} \ili{Turkish} fulfilled their role of state languages in several Empires, while regional groups still used their own varieties/languages and some vernacular of greater distribution. Language islands played a major role. Many languages of the \isi{minorities} were characterised by diglossia \citep{ferguson1959,ferguson1996}, the coexistence of daily spoken varieties of the language and a highly codified, formal written variety differing considerable in lexicon and grammar: Sephardic Jews in the \ili{Ottoman} Empire spoke \ili{Judeo-Spanish}, while \ili{Hebrew} constituted their sacred language. \ili{Aramaic} Christians in Syria under \ili{Ottoman} rule spoke varieties of \ili{Aramaic} and \ili{Arabic}, while classical \ili{Syriac} was used in the church. \isi{Literacy}, however, was not restricted to the high prestige variety, but the lower variety was written as well..

Only from the 19th century onwards did the concept of a national language, encompassing the entire population of a larger region, begin to spread from Europe to other parts of the world, including the \isi{MENA region}. In modern societies, proficiency\is{Language!proficiency} in the written standard language has become a prerequisite for all members of society. Over the past hundred years, literacy requirements have evolved from merely being able to sign a document or understand basic instructions to using written language for a wide range of purposes across work, education, and leisure. This literacy may encompass not only the official language of the host country but also other languages and varieties used within the migrant community or family.

The idea of language maintenance\is{Language!maintenance} \citep{fishman1995} itself is closely linked to transmitting the written variety of this language at school or in community education, leading to tensions between supporters of the ``good'' written language vs. the local spoken ``\isi{dialect}''. Otherwise, the official \isi{minority} language might be only in use in a small intellectual segment of the community, while the majority considers their spoken variety as a source of identity. The use of \isi{script} and even of orthographic features is closely linked to social processes of accommodation and divergence \citep{sebba2007,jaffe_androutsopoulos_sebba_johnson2012}.

It follows from the previous argument, that ``mother tongue'' in migration contexts is no longer acceptable as a scientific term. But ``first'' or ``family language'' and even ``language of origin'' poses some problems. In the last decades, the term ``heritage language'' was more frequently used. It describes ``… speakers as unbalanced bilinguals who acquired their first, heritage language in the home and whose second language serves as the dominant language of the broader community'' \citep[152]{scontras_putnam2020}. \isi{Heritage languages} do not differ from their original languages only by way of (uncontrolled) interference from the contact language used by the matrix society in the host country \citep{polinsky2011,polinsky2018}. Linguistic practices (including translanguaging, language learning strategies, etc.) are sometimes sustained for long periods of time in migration settings. The (structural, emblematic, semiotic) imprint this leaves on the resultant heritage language is not just an involuntary effect of language contact (and ``imperfect'' learning) but may be a strategy of appropriation and cultural adaptation. 

While heritage, including linguistic heritage, resonates with notions of stable tradition and transmission over long stretches of time. The condition of superdiversity evokes the idea of short-lived, ever-evolving, fragmented and fluid situations and phenomena. In such a way, the trajectories of languages brought to Germany still constitute an open project.

\section{Scope and purpose of the volume}

This edited volume is based on contributions from the conference ``Language, Society \& \isi{Identity} in Diverse Ethnolinguistic Contexts'', which was held in June 2023 in Frankfurt within the framework of the LOEWE research cluster.  The LOEWE\footnote{%
The acronym stands for: Landes-Offensive zur Entwicklung wissenschaftlich-ökonomischer Exzellenz (State initiative for the development of scientific and economic excellence)} research cluster ``Minority Studies: Language and \isi{Identity}'' (2020--2024) was an attempt to fill gaps in the field of \isi{minority} studies. Comprising researchers from Goethe University Frankfurt and University of Giessen it developed an interdisciplinary approach for studying problems of identity formation among \isi{minorities}. 

Aiming to address the five previously mentioned gaps in migration studies, this edited volume embarks on a comprehensive exploration of both historical and contemporary evolutions from the \isi{MENA region} to the \ili{German} \isi{diaspora}. By spanning across time, our approach ensures that the findings are not mere snapshots of the present. Instead, they are grounded in an appreciation of past trajectories, providing a richer context for understanding potential future developments. 

From a historical perspective, the focus is on Georgians in Georgia, Azerbaijan, and Iran and Sephardic Jews and speakers of \ili{Aramaic} varieties in the \ili{Ottoman} Empire. Regarding immigration in Germany, the studies include Kurds from Iraq and Syria, Pashtu speakers from Afghanistan, and Tarifit\il{Berber!Tarifit}-speaking Berbers from northern \isi{Morocco}.

While some studies concentrate on the linguistic, historical, and social aspects that shape the experiences of these \isi{minority} groups in their homelands, others address their experiences as migrants in Germany, a country of immigration.

The book delves into the nuanced processes of identity formation among \isi{minorities} from \isi{MENA region}, considering the cultural, social, and historical backgrounds that influence these processes. It emphasises the unique experiences of \isi{minority} groups that are marginalised within their own countries of origin, shedding light on how these internal dynamics influence their experiences and identities upon migrating. With respect to language, the tension between spoken vernaculars (\ili{Judeo-Spanish}, Kurmanji\il{Kurdish!Kurmanji} \ili{Kurdish}, several \ili{Aramaic} and \ili{Georgian} varieties, Tarifit\il{Berber!Tarifit} \ili{Berber}) and prestigious written languages (\ili{Hebrew}, \ili{Syriac}, \ili{Arabic}, \ili{Persian}, and \ili{Turkish}), not only in relation to official state languages in the countries where they reside but also within the communities themselves, is crucial. Even the choice of writing systems and scripts plays a central role in these processes.

One of the key objectives of this volume is to shift the focus of migration studies to include a more balanced view of both destination and origin countries. By doing so, it provides a deeper understanding of how migration processes are affected by the conditions in home countries. This includes an examination of how linguistic diversity, historical conflicts, and social stratification in these countries impact the identities of migrants.

Through a collection of essays and studies by various experts, the volume offers a multidimensional perspective on the challenges faced by these \isi{minority} groups. It aims to provide a more holistic understanding of their identity dynamics, contributing to more effective policies\is{Language!policies} and interventions for their integration and support in both their home and host countries. Therefore, three chapters deal with political and legal regulations in Germany and other European countries concerning linguistic and religious rights of the new \isi{minorities}.

This volume is an essential resource for scholars, policymakers, and anyone interested in the intricate relationship between migration, identity, language, and literacy, and the socio-political contexts of the MENA region. It serves as a crucial step towards acknowledging and addressing the complex realities of \isi{minorities} from this region.

\section{Overview of the volume}

The first part, ``Accommodation and interculturality in language acquisition\is{Language!acquisition} and use'', deals with the intricate dynamics of language transformation, adaptation, and cultural intermingling that result from migration and encounters with diverse cultures. It explores how these communities adjust their language use to blend into new environments and ``accommodate'' the language preferences of their culture and religion.

The paper \citetitle{chapters/1.1_IFLMD_Petrou}\is{discrimination}\il{Berber!Tarifit} by \citeauthor*{chapters/1.1_IFLMD_Petrou}
%``\isi{Discrimination} in communicative behaviours as perceived by 
% Tarifit\il{Berber!Tarifit} speakers'' by Petrou, Boutemin and Fanego Palat
provides a lens into the challenges faced by the speakers of Tarifit\il{Berber!Tarifit} \ili{Berber} in their communicative interactions. It further reinforces our understanding of the transformative and adaptive linguistic strategies these communities must employ to integrate into new cultural landscapes.

\citetitle{chapters/1.2_IFLMD_Topidi} by \citeauthor*{chapters/1.2_IFLMD_Topidi}
%``When words destroy worlds: Online hate speech against religious minorities %as harm'' by Kyriaki Topidi 
offers an exploration into the negative side of linguistic diversity and cultural encounter. It spotlights how language, particularly in online spaces, can be used as a tool for \isi{harm}, targeting \isi{minorities} including migrant groups. It enhances our understanding of how language use within migrant communities can be influenced by external factors, such as discrimination and \isi{hate speech}.

The second part, ``Heritage language education in Europe with a special focus on non-official migrant languages'', discusses the significance of heritage language education in Europe. It explores the opportunities and challenges faced by both immigrants and host societies in maintaining, teaching, and learning these heritage languages. It lays the groundwork by exploring the legal framework surrounding language rights and obligations for immigrant \isi{minorities}. It delves into how diversity governance can contribute to the integration of new \isi{minorities} into the societal fabric, while also respecting and preserving their unique cultural and linguistic heritage.

In Chapter 4, \citetitle{chapters/2.3_IFLMD_Hofmann}, \citeauthor*{chapters/2.3_IFLMD_Hofmann}
%Hofmann, Malkmus and Östermann 
offer an overview of the structure and legal framework of heritage language education programs in the \ili{German} state of \isi{Hesse}, as well as in selected regions across Europe. It explores the legislative and practical aspects that shape and define these programs, providing a comparative perspective across different geographical and cultural contexts. 

Chapter 5 by \citeauthor*{chapters/2.4_IFLMD_Medda-Windischer} \citetitle{chapters/2.4_IFLMD_Medda-Windischer}
%Medda-Windischer and Zeba 
delves into the language rights and obligations of new \isi{minority} groups, emerging from recent migration flows by leveraging an analysis of legal and policy documents, in conjunction with an exploration of prior literature and empirical studies undertaken in this area. The discussion provides a nuanced understanding of how language rights are perceived and implemented for these burgeoning \isi{minority} groups.

\citeauthor*{chapters/2.5_IFLMD_Koca}'s study, \citetitle{chapters/2.5_IFLMD_Koca} perspectives of newly arrived \ili{Kurdish} and \ili{Afghan} families in \ili{Germany} on their heritage languages and linguistic situation (Chapter 6), provides insightful perspectives on the dynamics of maintaining heritage languages amidst the process of learning official languages within migrant communities, juxtaposing the views of children and their parents at two different points of time.

The third part, ``Scripts as markers of identity'', delves into the historical dynamics of how writing systems, or scripts, have not only functioned as media for conveying messages, but as emblematic tokens of cultural and linguistic identity among diverse communities and migrant groups in the premodern \isi{MENA region}. It undertakes an exploration of how the phenomena of language contact and multilingualism, within the context of migration, have reshaped the evolution of written texts among diverse language communities. This includes the transformation of scripts, the borrowing and integration of words from one language to another, and even the birth of novel linguistic and cultural elements.

\citetitle{chapters/3.6_IFLMD_Mengozzi} (Chapter 7) \il{Syriac}\il{Neo-Aramaic}\il{Arabic!Garshuni} by \citeauthor*{chapters/3.6_IFLMD_Mengozzi} provides a compelling case study on how scripts serve as essential markers of identity. By tracing the evolution of \ili{Syriac}, \ili{Neo-Aramaic}, and Garshuni Arabic\il{Arabic!Garshuni} scripts, this paper explores how these scripts emerged and transformed over time, reflecting their communities' evolving cultural and linguistic identities.

\citeauthor*{chapters/3.7_IFLMD_Leibiusky}'s article \citetitle{chapters/3.7_IFLMD_Leibiusky}\is{Meam Loez} (Chapter 8) examines a specific instance of \isi{script} as an identity marker in the Jewish community of the \ili{Ottoman} Empire after the crisis of Sabbatai Zevi who claimed to be the long-awaited Messiah. Through the lens of Magriso's work, it investigates how scripts and language interweave to reflect and shape religious identity.

The article by \citeauthor*{chapters/3.8_IFLMD_Beridze} titled \citetitle{chapters/3.8_IFLMD_Beridze}\is{identity}\il{Georgian} (Chapter 9) offers a unique perspective on how scripts as well as features of spoken language can serve as emblems of identity within linguistically diverse communities. By focusing on \ili{Georgian} linguistic islands, the paper provides insights into the pivotal role of scripts in preserving and asserting cultural and linguistic identities.

\sloppy
\printbibliography[heading=subbibliography,notkeyword=this]
\end{document}
